% !PRINTED 613
\weberpage{613}
%<<CONT

4. \gesperrt{Die Anzahl der nicht äquivalenten Irrationalzahlen der Discriminante $\varDelta$ ist gleich der Anzahl der Formenclassen des Körpers $\Omega$}\footnote{Nach der am Schlusse des \S.\,128 des ersten Bandes gemachten Bemerkung über die \gesperrt{Gauss}'sche Theorie der quadratischen Formen stimmt die Classenzahl der quadratischen Formen nach der \gesperrt{Gauss}'schen Definition nicht immer überein mit der Anzahl der Idealclassen des Körpers $\Omega$. Um die Uebereinstimmung herzustellen, muss man in manchen Fällen die Idealclassen noch weiter in je zwei Classen theilen. Bei unserer Theorie der Aequivalenz der Zahlen $\omega$, wo die eigentliche nicht von der uneigentlichen Aequivalenz unterschieden wird, ist diese weitere Eintheilung nicht erforderlich. Vergl. \gesperrt{Dirichlet-Dedekind}, Zahlentheorie. 4. Auflage, S. 578, 584.}\gesperrt{.}

\begin{center}
§. 168.

\gesperrt{Composition der quadratischen Irrationalzahlen.}
\end{center}

Wir haben im \S.\,166 gesehen, dass zu jedem ganzen Functional $\varphi$ im Körper $\Omega$ eine gewisse quadratische Irrationalzahl $\omega$ von der Discriminante $\varDelta$ gefunden werden kann.

Wir müssen aber jetzt noch untersuchen, inwiefern diese Zahl $\omega$ durch den Körper $\Omega$ und durch das Functional $\varphi$ bestimmt ist.

Nach \S.\,166, (1) ist die Basis von $\varphi$

\weq{1}{\alpha_1 = a_{1,1},\quad \alpha_2 = a_{1,2}+a_{2,2}\Theta}

dadurch definirt, dass $a_{1,1}$ die kleinste positive ganze rationale Zahl ist, die durch $\varphi$ theilbar ist, und $a_{2,2}$ ist dann durch

\weq{2}{N_a(\varphi) = a_{1,1}a_{2,2}}

bestimmt, also sind die beiden Zahlen $a_{1,1},a_{2,2}$ durch $\varphi$ eindeutig bestimmt. Es fragt sich noch, inwiefern auch $a_{1,2}$ bestimmt ist. Die Zahl $a_{1,2}$ ist aber nach dem Modul $\varphi$ durch die Congruenz $a_{1,2}+a_{2,2}\Theta\equiv 0$ gegeben, und da $a_{1,2}$ auch eine ganze rationale Zahl sein muss, so ist $a_{1,2}$ bis auf ein Vielfaches von $a_{1,1}$ bestimmt. Dieses Vielfache von $a_{1,1}$ bleibt aber in der That willkürlich, da, wenn $\alpha_1,\alpha_2$ eine Basis von $\varphi$ ist, dasselbe für jedes ganze rationale $k$ auch von $\alpha_1,\alpha_2+k\alpha_1$ gilt.

Da nun $\alpha_2:\alpha_1=\omega$ die gesuchte quadratische Irrationalzahl ist, so erhalten wir den Satz:
%CONT>>
